\section{Hochschild homology  \texorpdfstring{$\BP\langle 2\rangle$}{BP2} with \texorpdfstring{$\BP\langle 1\rangle$}{BP1} coefficients}
This section is devoted to our computation of topological Hochschild homology of $\BP\langle 2\rangle$ with coefficients in $\BP\langle 1\rangle$. In Section~\ref{sec:SS}, we compute the spectral sequence \eqref{BrunSS} in the case $n=2$. In Section~\ref{sec:extensions}, we resolve extensions in this spectral sequence. 

\subsection{Computating the spectral sequence}\label{sec:SS}

We first use the known computation of $\THH_*(\BP\langle 2\rangle;H\mathbb{Z}_{(p)})$ to deduce differentials that will be useful for bootstrapping to the case of $\THH_*(\BP\langle 2\rangle;\BP\langle 1\rangle)$.

\begin{proposition}
  There is a spectral sequence
  \begin{equation}
    \THH_*(\BP\langle 1\rangle; H \Z_{(p)}) \langle \sigma v_2 \rangle \Rightarrow \THH_*(\BP \langle 2 \rangle; H \Z_{(p)}).
  \end{equation}

  The differentials in this spectral sequence are given for any $i \geq 1$ by
  \begin{equation}
    \begin{gathered}
      d_1(a_i) \dot{=} p^{\nu_p(i-1)} a_{i-1} \sigma v_2, \\
      d_1(b_i) \dot{=} p^{\nu_p(i-1)} b_{i-1} \sigma v_2.
    \end{gathered}
  \end{equation}
\end{proposition}

\begin{proof}
  The fiber sequence 
  \[ 
    \Sigma^{2p^2-1}\THH(\BP\langle 1\rangle)\to \THH(\BP\langle 2\rangle ;\BP\langle 1\rangle)\to \THH(\BP\langle 1\rangle)
  \]
  is of $\BP\langle 1\rangle \otimes_{\BP \langle 2 \rangle} \BP\langle 1\rangle^\op$-modules, so is compatible with the left $v_1$-multiplication map. This produce a fiber sequence
  \begin{equation}
    \Sigma^{2p^2-1}\THH(\BP\langle 1\rangle; H\Z_{(p)})\to \THH(\BP\langle 2\rangle ; H\Z_{(p)})\to \THH(\BP\langle 1\rangle; H\Z_{(p)})
  \end{equation}
  and a spectral sequence
  \begin{equation}
    \THH_*(\BP\langle 1\rangle; H \Z_{(p)}) \langle \sigma v_2 \rangle \Rightarrow \THH_*(\BP \langle 2 \rangle; H \Z_{(p)}).
  \end{equation}

  By \cite{ACH24}*{Theorem~3.18}, we know that 
  \[
    \THH_*(\BP\langle 1\rangle,\mathbb{Z}_{(p)})=\mathbb{Z}_{(p)}\langle \lambda_1\rangle \otimes (\mathbb{Z}_{(p)}\oplus T_0^1)
  \]
  and 
  \[
    \THH_*(\BP\langle 2\rangle,\mathbb{Z}_{(p)})=\mathbb{Z}_{(p)}\langle \lambda_1,\lambda_2\rangle \otimes (\mathbb{Z}_{(p)}\oplus T_0^2)
  \]
  So we know that there is long exact sequence 
  \[ 
    \to T_0^{2} \to T_0^1\overset{\partial}{\to} \Sigma^{2p^2}T_0^1 \to 
  \]
  where the connecting map $\partial$ correspond to the differential of the spectral sequence.
  Since 
  \[ 
    T_0^1=\bigoplus_{s\ge 1}\mathbb{Z}/p^s[(\mu^{p^2})^{p^s}]\{\lambda_{1+s}(\mu^{p^2})^{jp^{s-1}}: 0\le j\le p-2\}
  \]
  and 
  \[ 
    T_0^2=\bigoplus_{s\ge 1}\mathbb{Z}/p^s[(\mu^{p^3})^{p^s}]\{\lambda_{2+s}(\mu^{p^3})^{jp^{s-1}}: 0\le j\le p-2\}
  \]
  where 
  \[ 
    \lambda_{s}=\begin{cases} \lambda_s & \text{ if } 1\le s\le n+1 \\ 
      \lambda_{s-1}(\mu^{p^{n+1}})^{p^{s-(n+2)}(p-1)}
    \end{cases}
  \]
  
  The map $T_0^2 \rightarrow T_0^1$ has image $p \cdot T_0^1$, so that each generator in degree $*$ of $T_0^1$ must have image by $\partial$ the top of the multiplication-by-$p$-tower of the generator in degree $*-2p^2$ in $T_0^1$. This is exactly the formulas we claimed for the torsion. The differentials in the non-torsion part are zero for degree reasons.
\end{proof}

\begin{lemma} \label{lem:thhell1}
  In $\THH_*(\BP\langle 1\rangle)$, if $|\alpha| = |b_i|$ for some $i \geq 1$, and $\alpha \neq 0$, then $v_1^{p-1} \alpha \neq 0$; if moreover $v_1^{p+1} \alpha = 0$ then $v_1^p \alpha = 0$ and $\alpha \in \Z_{(p)} \{ v_0^{\nu_p(i)}b_i \}$. If $\alpha$ is a multiple of $v_1$, then $\alpha \in p\Z_{(p)} \{ v_0^{\nu_p(i)}b_i \}$ 
\end{lemma}

\begin{proof}
  Since $\THH_{|b_i|}(\BP\langle 1\rangle)$ for some $i\ge 1$ is generated as $\Z_{(p)}$-module by the element $v_1^k v_0^h b_j$ for some $h,\, k \geq 0$ and $j \geq 1$, it is sufficient to consider the case where $\alpha = v_1^k v_0^h b_j$. We consider the degrees:
  \begin{equation}
    \begin{gathered}
      |\alpha| = |b_i| \\
      \iff 2j p^2 - 1 + 2p - 1 + 2k(p-1) = 2i p^2 - 1 + 2p - 1 \\
      \iff (i - j)p^2 = k(p -1)
    \end{gathered}
  \end{equation}
  so that $p^2$ divides $k$; write $k = p^2 k'$. However, the order for multiplication by $v_1$ of $v_0^h b_j$ is given by the relation
  \begin{equation}
    v_1^{p^{\nu_p(j) - h + 1}+ \dots + p} v_0^h b_j = 0
  \end{equation}
  by Theorem~\ref{thm:AHL} so that since $\alpha \neq 0$,
  \begin{equation}
    k = p^2 k' \leq p^{\nu_p(j) - h + 1}+ \dots + p
  \end{equation}
  and thus
  \begin{equation}
    k + p - 1 \leq p^{\nu_p(j) - h + 1}+ \dots + p
  \end{equation}
  which means that $v_1^{p - 1} \alpha \neq 0$.

  If moreover $v_1^{p+1} \alpha = 0$, we have 
  \begin{equation}
    |v_1^{p+1} \alpha| = |v_1^{p+1}| + |b_i| = |b_{i+1}| - 2
  \end{equation}
  but the $v_1$-Bocksteins are supported by the classes $a_j$ of $\THH_*(\BP\langle 1\rangle; H\Z_{(p)})$ by \cite[Theorem~6.4]{AHL10}, whose degrees are of the form
  \begin{equation}
    |a_j| = |b_j| - 2p + 1.
  \end{equation}
  and thus it must be the case that $v_1^p \alpha = 0$.
  
  Finally, this means that the $v_1$-Bocksteins that kill $v_1^p \alpha$ and $v_1^p v_0^{\nu_p(i)} b_i$ are supported by some multiple of $p$ of the same class $a_{i + 1}$, so that $v_0^{\nu_p(i)} b_i$ is at the bottom of a tower of extensions that includes $\alpha$, and we conclude that  $\alpha \in \Z_{(p)}\{ v_0^{\nu_p(i)} b_i \}$. The last statement follows from the relations in Theorem~\ref{thm:AHL}. 
\end{proof}

\begin{lemma} \label{lem:thhell2}
  In $\THH_*(\BP\langle 1\rangle)$, if
  \begin{equation}
    |v_0^hb_j| \leq | b_i| < | v_1^{p^{\nu_p(j) - h + 1} + \dots + p} v_0^hb_j |
  \end{equation}
  for some $ 0 \leq j \leq i$ and $0 \leq h \leq \nu_p(j)$ then
  \begin{equation}
    | v_1^{p^{\nu_p(i) + 1} + \dots + p} b_i | \leq | v_1^{p^{\nu_p(j) - h + 1} + \dots + p} v_0^hb_j |.
  \end{equation}
\end{lemma}

\begin{proof}
  We have
  \begin{equation}
    \begin{aligned}
      | v_1^{p^{\nu_p(j) - h + 1} + \dots + p} v_0^hb_j |
      & = |b_j| + 2(p^{\nu_p(j) - h + 2} - p) \\
      & = |b_{j + p^{\nu_p(j) - h}} | - 2p
    \end{aligned}
  \end{equation}
  so that the hypothesis implies $j \leq i < j + p^{\nu_p(j) - h}$, and thus $\nu_p(i) < \nu_p(j) - h$. Considering the base-$p$ expansion of $i$ and $j$, we see that $ i + p^{\nu_p(i)} \leq j + p^{\nu_p(j) - h}$ which implies the result.
\end{proof}

\begin{proposition}\label{theorem:differentialellcoeff}
The differentials in the spectral sequence
\[ 
\THH_*(\BP\langle 1\rangle)\langle \sigma v_2\rangle \implies \THH_*(\BP\langle 2\rangle,\BP\langle 1\rangle)
\]
are given by 
\begin{equation}
  d_1(b_i)\dot{=} v_0^{\nu_p(i-1)}b_{i-1} \sigma v_2
\end{equation}
for each $i \geq 2$ and they are $v_1$-linear.
\end{proposition}

\begin{proof}
  First, we consider the rationalization of the fiber sequence 
  \[ 
    \Sigma^{2p^2-1}\THH(\BP\langle 1\rangle)\to \THH(\BP\langle 2\rangle ;\BP\langle 1\rangle)\to \THH(\BP\langle 1\rangle)
  \]
  This produces a spectral sequence 
  \[ 
    \THH_*(\BP\langle 1\rangle)\otimes \mathbb{Q}\langle \sigma v_2\rangle \implies \THH_*(BP\langle 2\rangle,\BP\langle 1\rangle)\otimes \mathbb{Q}
  \]
  which we know collapses by Proposition~\ref{prop:rational}. The map of spectral sequences 
  \[ 
    \begin{tikzcd}
      \THH_*(\BP\langle 1\rangle)\langle \sigma v_2\rangle \ar[d] \arrow[=>]{r} & \THH_*(BP\langle 2\rangle,\BP\langle 1\rangle)\ar[d] \\
      \THH_*(\BP\langle 1\rangle)\otimes \mathbb{Q}\langle \sigma v_2\rangle \arrow[=>]{r} & \THH_*(BP\langle 2\rangle,\BP\langle 1\rangle)\otimes \mathbb{Q}
    \end{tikzcd}
  \]
  is injective on the torsion free part and therefore the torsion free summands in the top spectral sequence consist of permanent cycles.

  The map of spectral sequences induced by $\BP\langle 1\rangle\to H\mathbb{Z}_{(p)}$
allows us to determine the differentials modulo $v_1$, so that
  \begin{equation}
    d_1(b_i) \dot{=} (v_0^{\nu_p(i-1)}b_{i-1} + \alpha ) \sigma v_2
  \end{equation}
  where $\alpha$ is a multiple of $v_1$. To prove our claim, it is sufficient to prove that $\alpha \in p \Z_{(p)}\{ v_0^{\nu_p(i-1)}b_{i-1} \}$, so that $v_0^{\nu_p(i-1)}b_{i-1} + \alpha \dot{=} v_0^{\nu_p(i-1)}b_{i-1}$ and thus the claimed formula  is true up to a different unit.

  Since
  \begin{equation}
    v_1^{p^{\nu_p(i)} + \dots + p } b_i = 0
  \end{equation}
  it must also be the case that
  \begin{equation} \label{eq:orderalpha}
    v_1^{p^{\nu_p(i)} + \dots + p } \alpha = 0.
  \end{equation}
  Let us write $k$ for the smallest integer such that $v_1^k \alpha = 0$. There are two possibilities:
  \begin{itemize}
  \item if $ |\alpha| < |b_i| \leq |v_1^k \alpha|$ then by Lemma~\ref{lem:thhell2} we have $|v_1^{p^{\nu_p(i)} + \dots + p } b_i| \leq |v_1^k \alpha|$. However, this implies that $|v_1^{p^{\nu_p(i)} + \dots + p } \alpha | < |v_1^k \alpha|$ so this is incompatible with equation \eqref{eq:orderalpha}.
  \item Otherwise, we have $|v_1^k \alpha | < |b_i|$. We know that $|\alpha | = |b_{i-1}|$ and that $|v_1^{p+1} b_{i-1}| < |b_i| < |v_1^{p+2} b_{i-1}|$ so that $k \leq p +1$. From Lemma~\ref{lem:thhell1}, we see that $v_1^p\alpha = 0$ and that $\alpha \in p\Z_{(p)}\{ v_0^{\nu_p(i-1)}b_{i-1} \}$.
  \end{itemize}
\end{proof}

\begin{remark}
The computation of the differential in Proposition~\ref{theorem:differentialellcoeff} can be considered as a computation of the action of the cap product by the class $e_{3}$ from Section~\ref{sec:bounding}, which lifts the class dual to $\mu_3\in \THH_{2p^3}(\BP\langle 2\rangle;\mathbb{F}_p)$. This cap product is computed explicitly on all classes modulo $v_1$ in~\cite[Corollary~5.7]{AHL10} and we combine this result with $v_1$-Bockstein spectral sequence information to produce the result. An explicit computation of the cap product on all of $\THH_*(\ell)$ does not appear in~\cite{AHL09} although they use certain formulas to produce hidden extensions. 
\end{remark}

We would now like to give a nice description of the kernel and cokernel of the $d_1$-differential. For this, we need the following lemma. 

\begin{lemma} \label{lem:thhell3}
  In $\THH_*(\BP\langle 1\rangle)$, $b_i$ and $v_0^{\nu_p(i-1)} b_{i-1}$ have the same order for multiplication by $p$ for any $ i \geq 2$ that is not a power of $p$. When $i = p^k$ for some $k \geq 1$, the order of $b_i$ is one more than that of $v_0^{\nu_p(i-1)} b_{i-1} = b_{i-1}$.
\end{lemma}

\begin{proof}
  If $p$ divides $i - 1$, then $p \cdot v_0^{\nu_p(i-1)} b_{i-1} = 0$ and $i \equiv 1 \pmod{p}$ so that $p \cdot b_i = 0$ too.

  If $p$ does not divide $i-1$ and $i$ is not a power of $p$, we can see from the description of $\THH_*(\BP\langle 1\rangle)$ in Theorem~\ref{thm:AHL} that the order $k$ for multiplication by $p$ of $v_0^{\nu_p(i-1)} b_{i-1} = b_{i -1}$ is such that $k -1$ is the number of digits that are $p - 1$ on the right of the base-$p$ writing of $i - 1$. Then $k-1$ is also the number of carry when adding one to $i -1$, and thus also the $p$-adic valuation of $i$, so that $k$ is the order for multiplication by $p$ of $b_i$.

  If $i$ is a power of $p$, there is one less extension in the tower above $b_{i-1}$, so that the order of $b_i$ is one more than that of $b_{i-1}$.
\end{proof}
We are now prepared to give our nicer description of the kernel and cokerel of the $d_1$-differential. 
\begin{proposition} \label{prop:specseqtorsion}
  In the spectral sequence
  \begin{equation}
    \THH_*(\BP\langle 1\rangle) \langle \sigma v_2 \rangle \Rightarrow \THH_*(\BP \langle 2 \rangle; \BP\langle 1\rangle)
  \end{equation}
  restricted to the torsion elements, the differential $d_1$ satisfies
  \begin{equation}
    \operatorname{coker} d_1 = \operatorname{coim}(\times v_1^p) \subset \THH_*(\BP\langle 1\rangle)\{ \sigma v_2 \}
  \end{equation}
  and
  \begin{equation}
    \ker d_1 =  \im (\times v_1^p) \oplus M \subset \THH_*(\BP\langle 1\rangle)\{ 1 \}
  \end{equation}
  where $\times v_1^p$ is the multiplication by $v_1^p$ map in the proper copy of $\THH_*(\BP\langle 1\rangle)$ and $M$ is the $\BP\langle 1\rangle_*$-submodule of $\THH_*(\BP\langle 1\rangle)$ generated by the $v_0^{k}b_{p^k}$ for all $k \geq 1$.
\end{proposition}

\begin{proof}
  The formula for the differentials of Proposition~\ref{theorem:differentialellcoeff} implies that $\im d_1 \subset \ker (\times v_1^p)$, since for any $i \geq 2$, $d_1(b_i) \in \Z_{(p)}\{ v_0^{\nu_p(i-1)} b_{i-1} \sigma v_2 \} \subset \ker (\times v_1^p)$. But Lemma~\ref{lem:thhell3} implies that
  \begin{equation}
    d_1(\BP\langle 1\rangle_*\{ b_i \} ) = \BP\langle 1\rangle_* \{ v_0^{\nu_p(i-1)} b_{i-1} \sigma v_2 \}
  \end{equation}
  and the description of $\THH_*(\BP\langle 1\rangle)$ implies that
  \begin{equation}
    \ker (\times v_1^p) = \bigoplus_{i \geq 2} \BP\langle 1\rangle_* \{ v_0^{\nu_p(i-1)} b_{i-1} \}
  \end{equation}
  so that
  \begin{equation}
    \im d_1 = \ker (\times v_1^p).
  \end{equation}

  On the other hand, the description of $\THH_*(\BP\langle 1\rangle)$ also implies that a class $\alpha \notin \im (\times v_1^p)$ is a sum of element of the form $v_1^k v_0^h b_i$ with $0 \leq k < p$. 
  But Lemma~\ref{lem:thhell3} implies that these elements, except for those in $M$, are exactly the sources of the non-zero differentials, so that
  \begin{equation}
    \ker d_1 = \im (\times v_1^p) \oplus M.
  \end{equation}
\end{proof}
